%%%%%%%%%%%%%%%%%%%%%%%%%%---Abstract---%%%%%%%%%%%%%%%%%%%%%%%%%%%%%%%%%%%
\newpage
\thispagestyle{plain}
\addcontentsline{toc}{chapter}{Abstract}

\begin{center}
\huge{\textbf{Abstract}}
\end{center}
\vspace{0.3cm}

Hybrid quantum--classical machine-learning models are increasingly proposed for
syndrome decoding and quantum error mitigation, but the evidence for them rests on
inconsistent baselines, quantum cost that is rarely normalized against accuracy, and
little testing outside the training distribution. This thesis benchmarks both tasks
under one protocol and organizes the findings around three questions: whether test
inputs are \emph{covered} by the training data, what each method \emph{costs}, and
whether its results are \emph{robust} to the split, seed, noise model and calibration.

In decoding, six decoders -- two matching, three classical learned and one hybrid
variational -- are compared on $16$ million repetition-code shots. Exact
minimum-weight perfect matching wins $33$ of $40$ configurations. Learned decoders
win only at distance $3$, where training covers $94.5\%$ of the syndrome space; at
distance $9$ it covers $1.65 \times 10^{-19}$, and they are $10$--$30\times$ worse.
The hybrid decoder beats its classical control on identical features at distance $3$
but not at distance $9$, and the neural decoder is about $740\times$ faster than a
pure-Python matching implementation.

Mitigation is posed as regression from noise-scaled expectation values to the ideal
value, on $4500$ circuit--observable pairs. Nine methods -- no mitigation, three
zero-noise-extrapolation variants, three classical regressors, a fidelity quantum
kernel and a data-reuploading variational circuit -- are compared on circuit-grouped
folds with paired bootstrap tests, over ten splits and ten seeds, under an idealized
Markovian model and a device-realistic model with verified non-Markovian memory
($\mathcal{N} = 0.88$, and $0$ under a reset control). In-distribution, the
three strongest learned methods (RMSE $0.0412$--$0.0436$) remove $41$--$44\%$ of the
error left by the best extrapolation ($0.0737$) in every split.

\emph{Coverage} decides out-of-distribution behaviour: on deeper circuits the neural
and quantum models are worse than no mitigation because depth and gate count leave
the training range, and without those inputs every learned method beats no
mitigation. \emph{Cost} separates models that tie on accuracy: the quantum kernel
ties a classical network ($0.0425 \pm 0.0012$ against $0.0426 \pm 0.0009$), but
estimated from shots it needs about $3 \times 10^8$ circuit executions per mitigated
value to approach that accuracy. \emph{Robustness} checks change the picture:
untuned bandwidth makes the kernel $47\%$ worse than no mitigation, the variational
model's error ranges from $0.056$ to $0.151$ across seeds, unitary folding halves the
kernel's margin over extrapolation ($38\%$ to $18\%$), memory worsens every
non-variational learned method but not extrapolation, and a $30\%$ recalibration
costs most learned methods no resolvable accuracy.

In both tasks, whether a learned method helps depends not on whether it is quantum,
but on coverage, cost, tuning and repeated initialisation. All results are simulated;
the decoding track uses the repetition code in place of the surface code and a single
training seed per decoder.

\vspace{0.4cm}
\noindent\textbf{Keywords:} syndrome decoding, minimum-weight perfect matching,
quantum error correction, quantum error mitigation, quantum machine learning,
zero-noise extrapolation, quantum kernel methods, variational quantum circuits,
non-Markovian noise, benchmarking.
